\begin{center}
    {\fontsize{12pt}{14.4pt}\selectfont\textbf{KATA PENGANTAR}\par}
    \vspace{0.5cm}
\end{center}

Puji syukur kami panjatkan ke hadirat Tuhan Yang Maha Esa atas rahmat dan karunia-Nya, sehingga kami dapat menyelesaikan penyusunan Laporan Magang dan Dokumentasi Sistem E-Learning ini dengan baik. Laporan ini disusun berdasarkan hasil pelaksanaan program magang kami di PT. BPR Jatim (Perseroda) pada Divisi Teknologi Informasi yang berlangsung dari tanggal 3 Februari sampai dengan 3 Juli 2026.

Dokumentasi ini dirancang sebagai panduan teknis operasional sistem e-learning yang dikembangkan, yang diharapkan dapat menjadi acuan bagi pihak kantor dalam menjalankan, memelihara, dan mengembangkan platform ini lebih lanjut.

Kami menyadari bahwa keberhasilan pelaksanaan magang dan penyusunan laporan ini tidak lepas dari bimbingan, bantuan, dan dukungan dari berbagai pihak. Oleh karena itu, kami ingin menyampaikan ucapan terima kasih yang sebesar-besarnya kepada:
\begin{enumerate}[label=\arabic*., leftmargin=*, noitemsep, topsep=6pt, parsep=0pt, itemsep=6pt]
    \item Pimpinan dan Manajemen PT. BPR Jatim (Perseroda) yang telah memberikan kesempatan kepada kami untuk belajar dan berkontribusi di perusahaan.
    \item Bapak Teezar Tioni selaku Kepala Divisi Teknologi Informasi PT. BPR Jatim (Perseroda) atas arahan, fasilitas, dan bimbingannya selama kami berada di bawah naungan Divisi TI.
    \item Bapak Achmad Fadli selaku Pembimbing Lapangan di Divisi TI yang senantiasa membimbing, mengarahkan, dan memberikan masukan teknis yang berharga dalam pengembangan sistem ini.
    \item Seluruh Staf dan Karyawan PT. BPR Jatim (Perseroda) yang telah bekerja sama dengan baik dan membantu kami dalam penyediaan data serta informasi yang dibutuhkan selama magang.
    \item Dosen Pembimbing dan segenap civitas akademika Politeknik Elektronika Negeri Surabaya (PENS) yang telah membimbing kami dari sisi akademis.
\end{enumerate}

Akhir kata, kami berharap dokumen panduan dan laporan akhir ini dapat memberikan manfaat yang nyata bagi operasional PT. BPR Jatim (Perseroda), khususnya dalam mendukung transformasi digital sistem pelatihan karyawan.

\vspace{0.8cm}
\begin{flushright}
Surabaya, 3 Juli 2026\\[1.6cm]
\textbf{Tim Penyusun}
\end{flushright}
\newpage
